%%%%%%%%%%%%%%%%%%%%%%%%%%%%%%%%%%%%%%%%%%%%%%%%%%%%%%%%%%%%%%%%%%%%%%%%%%%%%%%%
%%                          uAlberta Thesis Template                          %%
%%                                     by                                     %%
%%                               Daniel Aldrich                               %%
%%                               Version: 3.0.0                               %%
%%%%%%%%%%%%%%%%%%%%%%%%%%%%%%%%%%%%%%%%%%%%%%%%%%%%%%%%%%%%%%%%%%%%%%%%%%%%%%%%
%                                                                              %
%       Copyright (c) 2026 Daniel Aldrich                                      %
%                                                                              %
%       Permission is hereby granted, free of charge, to any person            %
%       obtaining a copy of this software and associated documentation         %
%       files (the "Software"), to deal in the Software without                %
%       restriction, including without limitation the rights to use,           %
%       copy, modify, merge, publish, distribute, sublicense, and/or           %
%       sell copies of the Software, and to permit persons to whom the         %
%       Software is furnished to do so, subject to the following               %
%       conditions:                                                            %
%                                                                              %
%       The above copyright notice and this permission notice shall be         %
%       included in all copies or substantial portions of the Software.        %
%                                                                              %
%       THE SOFTWARE IS PROVIDED "AS IS", WITHOUT WARRANTY OF ANY KIND,        %
%       EXPRESS OR IMPLIED, INCLUDING BUT NOT LIMITED TO THE WARRANTIES        %
%       OF MERCHANTABILITY, FITNESS FOR A PARTICULAR PURPOSE AND               %
%       NONINFRINGEMENT. IN NO EVENT SHALL THE AUTHORS OR COPYRIGHT            %
%       HOLDERS BE LIABLE FOR ANY CLAIM, DAMAGES OR OTHER LIABILITY,           %
%       WHETHER IN AN ACTION OF CONTRACT, TORT OR OTHERWISE, ARISING           %
%       FROM, OUT OF OR IN CONNECTION WITH THE SOFTWARE OR THE USE OR          %
%       OTHER DEALINGS IN THE SOFTWARE.                                        %
%                                                                              %
%%%%%%%%%%%%%%%%%%%%%%%%%%%%%%%%%%%%%%%%%%%%%%%%%%%%%%%%%%%%%%%%%%%%%%%%%%%%%%%%

% !TeX program = pdflatex
% !TeX spellcheck = en_CA
% !TeX TXS-program:bibliography = biber

% Write command to allow the nomenclature to be generated properly.
\immediate\write18{makeindex \jobname.nlo -s nomencl.ist -o \jobname.nls}
\documentclass[
	oneside,
	chapterbib,
	saychapapp,
	fancyheaders,
	colourlinks]{ualberta}
	% OPTIONS FOR ualberta.cls:
	% chapterbib - Automatically prints references at the end of each chapter.
	%
	% pdfa - To convert the PDF to PDF/A format (REQUIRED for GPS Submission)
	%
	% oneside - Standard for submitting to GPS.
	%
	% twoside - If you want to print your thesis double sided like a novel.
	%            Note: GPS requires the submission to be one sided. Please use
	%            the option `oneside' for submission to GPS.
	%
	% saychapapp - If you want your thesis to say Chapter # and Appendix @ in
	%               the ToC instead of just having the # or @.
	%               (GPS is inconsistent on if this is truly a Requirement.)
	%
	% fancyheader - If you want your thesis to have chapter headers rather than
	%               just a page number at the bottom.
	%
	% colourlinks - If you want your thesis to colour the crossref/urls.
	%               Remove if you want them to be all black.

% Option to change the Level of subheading included in the Table of Contents
%  This should be set at 2, 3, or 4 (As per GPS)
	\settoclevel{3}

%%%%%%%%%%%%%%%%%%%%%%%%%%%%%%%%%%%%%%%%%%%%%%%%%%%%%%%%%%%%%%%%%%%%%%%%%%%%%%%%
%                                FILE LOCATIONS                                %
%%%%%%%%%%%%%%%%%%%%%%%%%%%%%%%%%%%%%%%%%%%%%%%%%%%%%%%%%%%%%%%%%%%%%%%%%%%%%%%%
% The following Commands Can be set to change the lookup locations for files.
% LaTeX FILES LOCATION
	% .  - This folder
	% .. - Up one Folder
	\setlatexfiledir{../00_LaTeX_Files/}
	\insertlatexfile{includePackages} % LOAD ALL PACKAGES TO BE INCLUDED
	% !TeX root = ../guide
% EVERYTHING AFTER THIS POINT IS USED ONLY FOR FORMATTING THE EXAMPLES IN THIS
%  DOCUMENT (DO NOT USE THESE IN YOUR THESIS)
%  PLEASE DELETE THESE.
	\usepackage{float}

	\usepackage{fontawesome}
	\usepackage{lipsum}
	\usepackage{xcolor}

	\usepackage{showexpl}

	\usepackage{microtype}

	\geometry{%
		left =0.75in,
		right=1.00in,
	}

	\usepackage{enumitem}
	\setlist{noitemsep}

%% COLOUR DEFINITIONS
	\definecolor{signalRed}   {rgb}{0.631,0.149,0.176}
	\definecolor{signalOrange}{rgb}{0.796,0.380,0.098}
	\definecolor{signalYellow}{rgb}{0.910,0.749,0.157}
	\definecolor{signalBlue}  {rgb}{0.000,0.282,0.451}
	\definecolor{signalGreen} {rgb}{0.000,0.557,0.369}
	\definecolor{signalBlack} {rgb}{0.055,0.075,0.075}
	\definecolor{yellow}      {rgb}{0.8,0.8,0}
	\definecolor{grey}        {rgb}{0.3,0.3,0.3}
	\definecolor{gold}        {rgb}{0.4,0.4,0.0}
	\definecolor{str}         {rgb}{1,0,1}
	\definecolor{cmd}         {rgb}{0,0,1}
	\definecolor{env}         {rgb}{0.5,0,1}
	\definecolor{pkg}         {rgb}{0,0.6,0}
	\definecolor{opt}         {rgb}{1,0.5,0}


%% NEW COMMANDS
	\newcommand{\exampleText}{\lipsum[1][1-4]}
	\AtBeginDocument{\lstset{style=LaTeXStyle}}

	\makeatletter
		\newcommand\insertplotcode[1]{\@plotdir\detokenize{#1}}
	\makeatother

	\newcommand{\University}{University of Alberta}
	\newcommand{\Uni}{UofA}
	\newcommand{\Faculty}{Faculty of Graduate and Post-Doctorate Studies}
	\newcommand{\Fac}{GPS}

	\newcommand{\commenting}[1]{%
		\begin{singlespace}
			\noindent
			\begin{tabularx}{\textwidth}{@{}lX@{}}
				\textcolor{grey}{\faCommenting} & \textcolor{grey}{\textit{#1}}
			\end{tabularx}
		\end{singlespace}~\par}

	\newcommand{\danger}[1]{%
		\begin{singlespace}
			\noindent
			\begin{tabularx}{\textwidth}{@{}lX@{}}
				\color{white}\colorbox{signalRed}{\faWarning\ \textsc{\textbf{DANGER}}} & {\textbf{\MakeUppercase{#1}}}
			\end{tabularx}
		\end{singlespace}~\par}

	\newcommand{\warning}[1]{%
		\begin{singlespace}
			\noindent
			\begin{tabularx}{\textwidth}{@{}lX@{}}
				\color{signalBlack}\colorbox{signalOrange}{\faWarning\ \textsc{\textbf{WARNING}}} &  {\textbf{#1}}
			\end{tabularx}
		\end{singlespace}~\par}

	\newcommand{\caution}[1]{%
		\begin{singlespace}
			\noindent
			\begin{tabularx}{\textwidth}{@{}lX@{}}
				\color{signalBlack}\colorbox{signalYellow}{\faWarning\ \textsc{\textbf{CAUTION}}} & {#1}
			\end{tabularx}
		\end{singlespace}~\par}

	\newcommand{\notice}[1]{%
		\begin{singlespace}
			\noindent
			\begin{tabularx}{\textwidth}{@{}lX@{}}
				\color{white}\colorbox{signalBlue}{\textsc{\textbf{\textsl{NOTICE}}}} & {#1}
			\end{tabularx}
		\end{singlespace}~\par}

	\newcommand{\alert}[1]{%
		\begin{singlespace}
			\noindent
			\begin{tabularx}{\textwidth}{@{}lX@{}}
				\faWarning & {#1}
			\end{tabularx}
		\end{singlespace}~\par}

	\newcommand{\note}[1]{%
		\begin{singlespace}
			\noindent
			\begin{tabularx}{\textwidth}{@{}lX@{}}
				\textbf{\textsl{Note:}} & {\textit{#1}}
			\end{tabularx}
		\end{singlespace}~\par}

	\newcommand{\info}[1]{%
		\begin{singlespace}
			\noindent
			\begin{tabularx}{\textwidth}{@{}lX@{}}
				\textcolor{blue}{\faInfoCircle} & \textcolor{blue}{#1}
			\end{tabularx}
		\end{singlespace}~\par}

	\newcommand{\question}[1]{%
		\begin{singlespace}
			\noindent
			\begin{tabularx}{\textwidth}{@{}lX@{}}
				\textcolor{orange}{\faQuestionCircle} & \textbf{#1}
			\end{tabularx}
		\end{singlespace}~\par}

	\newcommand{\answer}[1]{%
		\begin{singlespace}
			\noindent
			\begin{tabularx}{\textwidth}{@{}lX@{}}
				\-\hspace{4ex}{\color{orange}{\faCheckCircle}} & \textit{#1}
			\end{tabularx}
		\end{singlespace}~\par}

	\newcommand{\qna}[2]{%
		\begin{singlespace}
			\noindent
			\begin{tabularx}{\textwidth}{@{}lX@{}}
				\textcolor{orange}{\faQuestionCircle} & \textbf{#1}
			\end{tabularx}
			\noindent
			\begin{tabularx}{\textwidth}{@{}lX@{}}
				\-\hspace{4ex} & \textit{#2}
			\end{tabularx}
		\end{singlespace}~\par}

	\newcommand{\tip}[1]{%
		\begin{singlespace}
			\noindent
			\begin{tabularx}{\textwidth}{@{}lX@{}}
				\textcolor{gold}{\faLightbulbO{}} & \textcolor{gold}{#1}
			\end{tabularx}
		\end{singlespace}~\par}

	\newcommand{\key}[1]{%
		\begin{singlespace}
			\noindent
			\begin{tabularx}{\textwidth}{@{}lX@{}}
				\textcolor{black}{\faKey} & {#1}
			\end{tabularx}
		\end{singlespace}~\par}

	\newcommand{\clouddownload}[2]{%
		\begin{singlespace}
			\noindent
			\begin{tabularx}{\textwidth}{@{}lX@{}}
				\faCloudDownload & \href{#1}{#2}
			\end{tabularx}
		\end{singlespace}~\par}

	\newcommand{\important}[1]{%
		\begin{singlespace}
			\noindent
			\begin{tabularx}{\textwidth}{@{}lX@{}}
				\textcolor{red}{\faExclamationCircle} & \textcolor{red}{#1}
			\end{tabularx}
		\end{singlespace}~\par}

	\newcommand{\tool}[1]{%
		\begin{singlespace}
			\noindent
			\begin{tabularx}{\textwidth}{@{}lX@{}}
				\textcolor{black}{\faWrench} & #1
			\end{tabularx}
		\end{singlespace}~\par}

	\newcommand{\setting}[1]{%
		\begin{singlespace}
			\noindent
			\begin{tabularx}{\textwidth}{@{}lX@{}}
				\textcolor{black}{\faCogs} & #1
			\end{tabularx}
		\end{singlespace}~\par}

	\newcommand{\software}[1]{%
		\begin{singlespace}
			\noindent
			\begin{tabularx}{\textwidth}{@{}lX@{}}
				\textcolor{black}{\faFileCodeO} & #1
			\end{tabularx}
		\end{singlespace}~\par}

	\newcommand{\option}[1]{%
		\begin{singlespace}
			\noindent
			\begin{tabularx}{\linewidth}{@{}X@{}}
				`\textbf{#1}'
			\end{tabularx}
		\end{singlespace}~\par}

	\newcommand{\file}[1]{%
		\begin{singlespace}
			\noindent
			\begin{tabularx}{\linewidth}{@{}X@{}}
				\faFileO{} \texttt{#1}
			\end{tabularx}
		\end{singlespace}~\par}

	\newcommand{\folder}[1]{%
		\begin{singlespace}
			\noindent
			\begin{tabularx}{\linewidth}{@{}X@{}}
				\faFolderO{} \texttt{#1}
			\end{tabularx}
		\end{singlespace}~\par}

	\newcommand{\str}[1]{%
		\texttt{\textcolor{str}{#1}}}

	\newcommand{\pkg}[1]{%
		\texttt{\textcolor{pkg}{#1}}}

	\newcommand{\cmd}[1]{%
		\texttt{\textcolor{cmd}{\textbackslash\detokenize{#1}}}}

	\newcommand{\opt}[1]{%
		\texttt{\textcolor{opt}{\detokenize{#1}}}}
	\newcommand{\oopt}[1]{%
		\texttt{[\textcolor{opt}{\detokenize{#1}}]}}
	\newcommand{\mopt}[1]{%
		\texttt{\{\textcolor{opt}{\detokenize{#1}}\}}}

	\newcommand{\env}[1]{%
		\texttt{\textcolor{env}{#1}}}

	\newcommand{\benv}[1]{%
		\texttt{\cmd{begin}\{\textcolor{env}{#1}\}}}
	\newcommand{\eenv}[1]{%
		\texttt{\cmd{end}\{\textcolor{env}{#1}\}}} % LOAD PACKAGES/MACROS FOR THE GUIDE

% PREFATORY LOCATION
	% .  - This folder
	% .. - Up one Folder
	\setprefatorydir{./01_Prefatory/}

% CHAPTER LOCATION
	% .  - This folder
	% .. - Up one Folder
	\setchapterdir{./02_Chapters/}

% BIBLIOGRAPHY LOCATION
	% NOTE: if you add bibliography entries after a compilation, you might notice
	%  references marked `[0]' to fix this just delete the auxiliary files.
	%  (*.aux, *.bbl, ... etc)
	%
	% .  - This folder
	% .. - Up one Folder
	\addbibresource{./03_References/References.bib}

% APPENDICES LOCATION
	% .  - This folder
	% .. - Up one Folder
	\setappendixdir{./04_Appendices/}

% MEDIA LOCATIONS
	% .  - This folder
	% .. - Up one Folder
	\setmediadir{./99_Inclusions/}
	\setimagedir{Images/}
	\settabledir{Tables/}
	\setplotdir{Plots/}
	\setcodedir{Code/}
	\setdatadir{Data/}
	\setpdfdir{PDFs/}
%%%%%%%%%%%%%%%%%%%%%%%%%%%%%%%%%%%%%%%%%%%%%%%%%%%%%%%%%%%%%%%%%%%%%%%%%%%%%%%%
%                         ADD ADDITIONAL RESOURCES                             %
%%%%%%%%%%%%%%%%%%%%%%%%%%%%%%%%%%%%%%%%%%%%%%%%%%%%%%%%%%%%%%%%%%%%%%%%%%%%%%%%
	\insertlatexfile{includeTheorems}
	\insertlatexfile{includeMacros}
	\insertlatexfile{listingCodeFormatting}

%%%%%%%%%%%%%%%%%%%%%%%%%%%%%%%%%%%%%%%%%%%%%%%%%%%%%%%%%%%%%%%%%%%%%%%%%%%%%%%%
%                  TITLE PAGE AND FRONTMATTER INFORMATION                      %
%%%%%%%%%%%%%%%%%%%%%%%%%%%%%%%%%%%%%%%%%%%%%%%%%%%%%%%%%%%%%%%%%%%%%%%%%%%%%%%%
% TITLE PAGE INFO
	\title{uAlberta Thesis User/Writing Guide}   % Title of your Thesis
	\author{Daniel R.\ Aldrich}                  % Your Full Name
	\degree{\insertlatexfile{selectDegree}}      % Uncomment Degree in file
	\specialization{}                            % Leave blank if none
	\deptfac{\insertlatexfile{selectDepartment}} % Uncomment Department in file
	\convocationdate{\the\year}                  % Convocation Year

% SELECT IF YOUR THESIS IS TO BE PUBLISH UNDER CC
	\creativecommons{\insertlatexfile{selectCreativeCommons}}

% AN EASY WAY TO INCLUDE OR EXCLUDE THE FOLLOWING SECTIONS IS TO COMMENT OR 
%  UNCOMMENT THE FOLLOWING LINES.
	% ABSTRACT
		\setabstracttext{}
	
	% PREFACE
		\setprefacetext{}
	
	% DEDICATION
		\setdedicationtext{}
	
	% QUOTE
		\setquotetext{}
	
	% ACKNOWLEDGEMENTS
		\setacknowledgementtext{}


%%%%%%%%%%%%%%%%%%%%%%%%%%%%%%%%%%%%%%%%%%%%%%%%%%%%%%%%%%%%%%%%%%%%%%%%%%%%%%%%
%                  NOMENCLATURE, GLOSSARY, ACRONYMS, ETC                       %
%%%%%%%%%%%%%%%%%%%%%%%%%%%%%%%%%%%%%%%%%%%%%%%%%%%%%%%%%%%%%%%%%%%%%%%%%%%%%%%%

% NOMENCLATURE
	\insertprefatory{Nomenclature}

% ACRONYMS
	\insertprefatory{Acronyms}

% GLOSSARY
	\insertprefatory{Glossary}


%%%%%%%%%%%%%%%%%%%%%%%%%%%%%%%%%%%%%%%%%%%%%%%%%%%%%%%%%%%%%%%%%%%%%%%%%%%%%%%%
%                              BEGIN DOCUMENT                                  %
%%%%%%%%%%%%%%%%%%%%%%%%%%%%%%%%%%%%%%%%%%%%%%%%%%%%%%%%%%%%%%%%%%%%%%%%%%%%%%%%
\begin{document}

% These are ``Magic'' Macros... no need to comment them out. If their sections 
%  are empty then they will automagically not be included.
	\maketitle                   % Creates the Title Page
	\makeabstract                % Creates the Abstract Page
	\makepreface                 % Creates the Preface Page
	\makededicationandquote      % Creates the Quote/Dedication Page
	\makeacknowledgements        % Creates the Acknowledgements Page

% SET ToC...etc SPACING
	%\singlespacing              % 1.00x Spacing
	\onehalfspacing              % 1.50x Spacing
	%\doublespacing              % 2.00x Spacing

% ALL LISTS/TABLES OF ________
% Everything below should Automatically be included if it has content and
%  excluded if it does not. There is no need to comment any line here. If you
%  would like a different order however feel free to move items to new lines.
	\tableofcontents             % Creates the Table of Contents
	\listoftables                % Creates the List of Tables
	\listoffigures               % Creates the List of Figures
	\listofplates                % Creates the List of Plates (photographs)
	\listofillustrations         % Creates the List of Illustration
	\listofvideos                % Creates the List of Videos
	\listofsounds                % Creates the List of Sounds
	\listofphotos                % Creates the List of Photos
	\listofsymbols               % Creates the List of Symbols (Nomenclature)
	\abbreviations               % Creates the List of Acronyms (Abbreviations)
	\glsaddall                   % Required for List of Acronyms and Glossary
	\generateglossary            % Creates the Glossary of Terms


%%%%%%%%%%%%%%%%%%%%%%%%%%%%%%%%%%%%%%%%%%%%%%%%%%%%%%%%%%%%%%%%%%%%%%%%%%%%%%%%
%                           ADD YOUR CONTENT HERE                              %
%%%%%%%%%%%%%%%%%%%%%%%%%%%%%%%%%%%%%%%%%%%%%%%%%%%%%%%%%%%%%%%%%%%%%%%%%%%%%%%%
	\bodyoftext % Switches the style of the document to that of the body

% SET DOCUMENT SPACING
	%\onehalfspacing       % 1.50x Spacing **Minimum Spacing for the THESIS BODY
	\doublespacing         % 2.00x Spacing
	%\triplespacing        % 3.00x Spacing

% To insert chapters from a separate tex file use the following commands
%  \insertchapter automatically looks in the Chapters folder and executes some
%  additional commands for generating end of Chapter References
%
% \input is the standard way of including a separate tex file where the 
%  \printbibliography command will generate the end of Chapter references.
% \input{./02_Chapters/ExampleChapter}
% \begin{singlespace}
%       \printbibliography[%
%               segment=\therefsegment,
%               heading=subbibintoc,
%               title=References]
% \end{singlespace}

	\insertchapter{01_Introduction}

\part{Writing a Thesis at the \University}\label{prt:writing}
	\insertchapter{XX_Requirements}
	\insertchapter{XX_Submitting_Your_Thesis}

\part{Writing with uAlberta Thesis \LaTeX\ Template}\label{prt:writingwithlatex}
	The following part of the guide is focused on the practical application of \LaTeX.
	Further, I will be showing you how to \LaTeX\ within the scope of the provided template. 
	Because, while learning \LaTeX\ provides you with a powerful typesetting tool, the template acts as the \enquote{rulebook} that ensures your document will automatically adhere to the  formatting guidelines set by the \Faculty\ (\Fac). 

	Using the template with \LaTeX\ effectively means you can spend your time on your research and data rather than tweaking margins, title page layout, or the ordering of front matter and other components of your thesis. 
	These sections will walk you through the basic use of \LaTeX\ and how the template, folder/file structure, and pre-defined macros/commands make the process of writing your thesis more streamlined.

	I will be covering the workflow from the ground up:
	\begin{itemize}
			\item \textbf{Setup:} Installing \LaTeX\ and the initialization of the template and understanding the \enquote{Main} document file (\path{ualberta.tex}).
			\item \textbf{Structure:} Organizing your chapters, abstracts, and inclusions, as well as utilizing the template's commands for automatically generating the thesis' front matter (Abstract, Preface, and Acknowledgements).
			\item \textbf{Content:} Best practices for inserting figures, tables, and complex equations.
			\item \textbf{Visuals:} Incorporating high-quality plots and graphs that maintain consistency in the final PDF.
			\item \textbf{Citations:} Managing a robust bibliography that integrates seamlessly with \University's required citation requirements.
	\end{itemize}

	\insertchapter{02_Getting_Started}
	\insertchapter{03_Document_Structure}
	\insertchapter{04_Figures_Tables}
	\insertchapter{05_Plots_And_Graphs}
	\insertchapter{06_Mathematical_Equations}
	\insertchapter{07_Citations_And_References}

% How to use the Recommended Softwares
%% Technical Writing and Bibliography
\part{Technical Writing and Reference Management}\label{prt:technicalwriting}
	\insertchapter{08_TeXstudio}
	
	\insertchapter{08_JabRef}

	\insertchapter{08_Notepad++}
\part{Data Analysis and Programming}\label{prt:dataandprog}
%% Data Analysis and Programming
	\insertchapter{11_GNU_Octave}

	\insertchapter{13_GitHub}

	\insertchapter{14_Calcpad}

\part{Images and Technical Figures}\label{prt:imgsandfigs}
%% Images and Technical Figures
	\insertchapter{15_ImageJ}

	\insertchapter{16_GIMP}

	\insertchapter{17_Dia}

%%%%%%%%%%%%%%%%%%%%%%%%%%%%%%%%%%%%%%%%%%%%%%%%%%%%%%%%%%%%%%%%%%%%%%%%%%%%%%%%
%                                 BIBLIOGRAPHY                                 %
%%%%%%%%%%%%%%%%%%%%%%%%%%%%%%%%%%%%%%%%%%%%%%%%%%%%%%%%%%%%%%%%%%%%%%%%%%%%%%%%

% Uncomment the \nocite{*} if you want to include works read, but not cited
%   within the body of your work.
	% \nocite{*}
	\singlespacing
	\printbibliography[heading=bibintoc]


%%%%%%%%%%%%%%%%%%%%%%%%%%%%%%%%%%%%%%%%%%%%%%%%%%%%%%%%%%%%%%%%%%%%%%%%%%%%%%%%
%                                  APPENDICES                                  %
%%%%%%%%%%%%%%%%%%%%%%%%%%%%%%%%%%%%%%%%%%%%%%%%%%%%%%%%%%%%%%%%%%%%%%%%%%%%%%%%
\appendix
% SET APPENDIX SPACING
	%\onehalfspacing              % 1.50x Spacing
	\doublespacing                % 2.00x Spacing
	%\triplespacing               % 3.00x Spacing

% To insert appendices from a separate tex file use the following commands
% \insertappendix automatically looks in the Appendices folder
%
% \input is the standard way of including a separate tex file
% \input{"./Appendices/PDF_Appendix.tex"}

	\insertappendix{0A_Additional_Figures}

	\insertappendix{0B_Additional_Tables}

	\insertappendix{0C_Code_Listings}

	\insertappendix{0D_Including_PDFs}

	\insertappendix{0X_Math_Lettering}

\end{document}
